% !TEX root = main.tex
% ============================================================
% 原文档导言区与图宏移植
% ============================================================
\usepackage{graphicx}
\usepackage{booktabs}
\usepackage{longtable}
\usepackage{array}
\usepackage{tabularx}
\usepackage{enumitem}
\usepackage{xcolor}
\usepackage{float}
\usepackage{needspace}
\usepackage{pdflscape}
\usepackage{hyperref}

% 表格正文列一律左对齐：中文两端对齐会把字距拉开，长串枚举时尤其明显
\newcolumntype{L}[1]{>{\raggedright\arraybackslash}p{#1}}
\newcolumntype{Y}{>{\raggedright\arraybackslash}X}

\hypersetup{colorlinks=true, linkcolor=black, urlcolor=blue, citecolor=black,
  pdftitle={系统分析与设计课程作业},
  pdfsubject={玉器门店顾客服务与经营管理系统：愿景、业务建模、需求、分析与设计},
  pdfauthor={张昊轩}}

% 图题与表题：标签与文字之间用 quad 分隔
\captionsetup{labelsep=quad}

% 图表按章编号，形如 图1-1、表4-12
\counterwithin{table}{chapter}
\counterwithin{figure}{chapter}
\renewcommand{\thetable}{\thechapter-\arabic{table}}
\renewcommand{\thefigure}{\thechapter-\arabic{figure}}

\newcommand{\codefont}{\ttfamily\fontsize{10}{11}\selectfont}
\newcommand{\code}[1]{\codefont #1}

% 下划线处允许折行，长字段名（如 transfer_apply_id）不至于压到右侧一列
\let\csuUnderscore\_
\renewcommand{\_}{\csuUnderscore\allowbreak}

% 带圈数字划归 CJK 类，由中文字体渲染，避免静默丢字
\xeCJKDeclareCharClass{CJK}{"2460 -> "24FF, "2014}

% ------------------------------------------------------------
% 图宏：依次在 ../figs/ 的分类子目录里找文件，都没有时给出提示方框
% ------------------------------------------------------------
\newlength{\figw}
\newlength{\figmaxh}
\setlength{\figmaxh}{22.5cm}

\newcommand{\fig}[3][1.0]{%
  \setlength{\figw}{#1\textwidth}%
  \clearpage
  \begin{figure}[H]
    \centering
    \IfFileExists{../figs/analysis/#3.png}{%
      \includegraphics[width=\figw,height=\figmaxh,keepaspectratio]{../figs/analysis/#3.png}%
    }{%
    \IfFileExists{../figs/analysis/#3}{%
      \includegraphics[width=\figw,height=\figmaxh,keepaspectratio]{../figs/analysis/#3}%
    }{%
    \IfFileExists{../figs/design/#3.png}{%
      \includegraphics[width=\figw,height=\figmaxh,keepaspectratio]{../figs/design/#3.png}%
    }{%
    \IfFileExists{../figs/design/#3}{%
      \includegraphics[width=\figw,height=\figmaxh,keepaspectratio]{../figs/design/#3}%
    }{%
      \fbox{\parbox[c][4.6cm][c]{0.8\textwidth}{\centering
        【待插入图形】\vspace{6pt}
        \texttt{../figs/#3}}}}}}}
    \caption{#2}
    \label{fig:#3}
  \end{figure}%
}

\newcommand{\figtwo}[5][1.0]{%
  \setlength{\figw}{#1\textwidth}%
  \clearpage
  \begin{figure}[H]
    \centering
    \IfFileExists{../figs/analysis/#3.png}{\includegraphics[width=\figw]{../figs/analysis/#3.png}}{%
    \IfFileExists{../figs/design/#3.png}{\includegraphics[width=\figw]{../figs/design/#3.png}}{\fbox{【待插入图形】\texttt{#3}}}}%
    \caption{#2}\label{fig:#3}\vspace{6pt}
    \IfFileExists{../figs/analysis/#5.png}{\includegraphics[width=\figw]{../figs/analysis/#5.png}}{%
    \IfFileExists{../figs/design/#5.png}{\includegraphics[width=\figw]{../figs/design/#5.png}}{\fbox{【待插入图形】\texttt{#5}}}}%
    \caption{#4}\label{fig:#5}
  \end{figure}%
}

% 同页两图、各自独立宽度：#1=宽1 #2=图题1 #3=图名1 #4=宽2 #5=图题2 #6=图名2
\newcommand{\figtwoW}[6]{%
  \clearpage
  \begin{figure}[H]
    \centering
    \IfFileExists{../figs/analysis/#3.png}{\includegraphics[width=#1\textwidth]{../figs/analysis/#3.png}}{%
    \IfFileExists{../figs/design/#3.png}{\includegraphics[width=#1\textwidth]{../figs/design/#3.png}}{\fbox{【待插入图形】\texttt{#3}}}}%
    \caption{#2}\label{fig:#3}\vspace{6pt}
    \IfFileExists{../figs/analysis/#6.png}{\includegraphics[width=#4\textwidth]{../figs/analysis/#6.png}}{%
    \IfFileExists{../figs/design/#6.png}{\includegraphics[width=#4\textwidth]{../figs/design/#6.png}}{\fbox{【待插入图形】\texttt{#6}}}}%
    \caption{#5}\label{fig:#6}
  \end{figure}%
}

\newcommand{\figab}[3][0.86]{%
  \needspace{10\baselineskip}
  \begin{figure}[H]
    \centering
    \IfFileExists{../figs/business-use-cases/#3}{%
      \includegraphics[width=#1\textwidth,height=\figmaxh,keepaspectratio]{../figs/business-use-cases/#3}%
    }{%
    \IfFileExists{../figs/business-sequence/#3}{%
      \includegraphics[width=#1\textwidth,height=\figmaxh,keepaspectratio]{../figs/business-sequence/#3}%
    }{%
    \IfFileExists{../figs/improved-business-sequence/#3}{%
      \includegraphics[width=#1\textwidth,height=\figmaxh,keepaspectratio]{../figs/improved-business-sequence/#3}%
    }{%
    \IfFileExists{../figs/system-use-cases/#3}{%
      \includegraphics[width=#1\textwidth,height=\figmaxh,keepaspectratio]{../figs/system-use-cases/#3}%
    }{%
      \fbox{\parbox[c][4.6cm][c]{0.8\textwidth}{\centering
        【待插入图形】\\[6pt]
        \texttt{../figs/#3}}}}}}}
    \caption{#2}
    \label{fig:#3}
  \end{figure}%
}

% ------------------------------------------------------------
% 整页图：图与图注在版心内水平、垂直双向居中，独占一页，不随正文浮动
% 用于本来就占一页的图，免得图贴页顶、下方空出一大片
% ------------------------------------------------------------
\newcommand{\figp}[3][1.0]{%
  \clearpage
  \setlength{\figw}{#1\textwidth}%
  \noindent
  \begin{minipage}[c][\textheight][c]{\textwidth}%
    \centering
    \IfFileExists{../figs/analysis/#3.png}{%
        \includegraphics[width=\figw,height=\figmaxh,keepaspectratio]{../figs/analysis/#3.png}%
      }{%
      \IfFileExists{../figs/design/#3.png}{%
          \includegraphics[width=\figw,height=\figmaxh,keepaspectratio]{../figs/design/#3.png}%
        }{%
        \fbox{\parbox[c][4.6cm][c]{0.8\textwidth}{\centering 【待插入图形】\ \texttt{../figs/#3}}}%
      }%
    }%
    \par\vspace{6pt}%
    \captionof{figure}{#2}\label{fig:#3}%
  \end{minipage}%
  \clearpage
}

\newcommand{\figpab}[3][0.86]{%
  \clearpage
  \noindent
  \begin{minipage}[c][\textheight][c]{\textwidth}%
    \centering
    \IfFileExists{../figs/business-use-cases/#3}{%
        \includegraphics[width=#1\textwidth,height=\figmaxh,keepaspectratio]{../figs/business-use-cases/#3}%
      }{%
      \IfFileExists{../figs/business-sequence/#3}{%
          \includegraphics[width=#1\textwidth,height=\figmaxh,keepaspectratio]{../figs/business-sequence/#3}%
        }{%
        \IfFileExists{../figs/improved-business-sequence/#3}{%
            \includegraphics[width=#1\textwidth,height=\figmaxh,keepaspectratio]{../figs/improved-business-sequence/#3}%
          }{%
          \IfFileExists{../figs/system-use-cases/#3}{%
              \includegraphics[width=#1\textwidth,height=\figmaxh,keepaspectratio]{../figs/system-use-cases/#3}%
            }{%
            \fbox{\parbox[c][4.6cm][c]{0.8\textwidth}{\centering 【待插入图形】\ \texttt{../figs/#3}}}%
          }%
        }%
      }%
    }%
    \par\vspace{6pt}%
    \captionof{figure}{#2}\label{fig:#3}%
  \end{minipage}%
  \clearpage
}
